\section{A vector representation that certifies tree packings}
\label{sec:auxiliary}

The next representation gives a matrix inequality sufficient for all the
partition inequalities in Theorem~\ref{thm:packing}. Its construction uses
strictly positive edge weights, which preserve every subset rank.
The normalization uses the log-determinant variational method behind
radial isotropic position~\cite[Section~2.2]{Barthe1998}. Here the target
vector of edge weights is the average of the edge-set indicator vectors
of the given disjoint trees. We give the approximate positive-weight form
that keeps every subset rank unchanged.

\begin{lemma}[Isotropic representation preserving ranks]
\label{lem:auxiliary-representation}
Let $B$ be the union of $q\geq1$ edge-disjoint spanning trees on an
$n$-vertex set, where $n\geq2$.  For every $\rho>0$ there are vectors
$z_e\in\mathbb R^{n-1}$, indexed by the individual edges of $B$, such that
\begin{equation}
 \sum_{e\in B}z_ez_e^{\mathsf T}=I_{n-1},
 \qquad \|z_e\|^2\leq\frac{1+\rho}{q}.
 \label{eq:auxiliary-isotropy}
\end{equation}
Moreover, there are an invertible linear map $M:\mathbf1^\perp\to
\mathbb R^{n-1}$ and positive numbers $s_e$ such that $z_e=s_eMb_e$.
Consequently the incidence vectors and the $z_e$ have the same rank on
every subset of edges.
\end{lemma}

\begin{proof}
Choose orthonormal coordinates on $\mathbf1^\perp$ and write $b_e$ for
the incidence columns in these coordinates.  Denote the given trees by
$T_1,\ldots,T_q$.  For $a\in\mathbb R^{E(B)}$, set
\[
 N(a)=\sum_{e\in B}\exp(a_e)b_eb_e^{\mathsf T},
 \qquad
 h(a)=\log\det N(a)-\frac1q\sum_{e\in B}a_e.
\]
Every $N(a)$ is positive definite: the columns of $T_1$ are a basis and
all its weights are positive.  Thus $h$ is a smooth, everywhere finite
function.

Let $c_t>0$ be the square of the determinant of the basis of columns
indexed by $T_t$.  Cauchy--Binet, followed by the arithmetic--geometric
mean inequality, gives
\begin{align*}
 \det N(a)
 &\geq\sum_{t=1}^q c_t\exp\left(\sum_{e\in T_t}a_e\right)\\
 &\geq q\left(\prod_{t=1}^q c_t\right)^{1/q}
       \exp\left(\frac1q\sum_{e\in B}a_e\right).
\end{align*}
Here the last exponent uses the fact that the $T_t$ partition $E(B)$.
In particular,
\[
 h(a)\geq L:=\log q+\frac1q\sum_{t=1}^q\log c_t.
\]
For $\gamma>0$, the function $h(a)+\gamma\|a\|^2$ is continuous and
coercive, hence has a minimizer $a^\gamma\in\mathbb R^{E(B)}$.
Comparison with $a=0$ shows that, with $K=h(0)-L\geq0$,
\[
 \gamma\|a^\gamma\|^2\leq K,
 \qquad
 \|\nabla h(a^\gamma)\|
 =2\gamma\|a^\gamma\|\leq2\sqrt{\gamma K}.
\]
The stationarity equation supplies the equality in the second formula.
Thus one may choose a finite $a=a^\gamma$ for which every coordinate of
$\nabla h(a)$ has absolute value at most $\rho/q$.  This remains valid
when $K=0$, since the displayed gradient bound is then zero.

The determinant derivative formula yields
\[
 \frac{\partial h}{\partial a_e}
 =\exp(a_e)b_e^{\mathsf T}N(a)^{-1}b_e-\frac1q.
\]
Define $z_e=\exp(a_e/2)N(a)^{-1/2}b_e$.  Their outer products sum
exactly to the identity, and the gradient bound gives their asserted
norm bound.  The common map $N(a)^{-1/2}$ is invertible and every
$\exp(a_e/2)$ is positive.  These operations preserve all subset ranks,
including the dependencies among parallel edges.
\end{proof}

\begin{lemma}[Projected trace and packing]
\label{lem:projected-trace-packing}
With the vectors in Lemma~\ref{lem:auxiliary-representation}, suppose
$A\subseteq E(B)$ and $c>0$ satisfy
\[
 \sum_{e\in A}z_ez_e^{\mathsf T}\succeq cI_{n-1}.
\]
Then $(V,A)$ contains at least
$\lfloor cq/(1+\rho)\rfloor$ edge-disjoint spanning trees.
\end{lemma}

\begin{proof}
Fix a partition $\mathcal P=\{V_1,\ldots,V_d\}$ into nonempty parts.
Let $W$ be the span of the $z_e$ indexed by all edges of $B$ whose
endpoints lie in one part.  The corresponding incidence vectors lie in
the direct sum of the zero-sum spaces on the parts.  The representation
therefore gives
\[
 \dim W\leq\sum_{j=1}^d(|V_j|-1)=n-d.
\]
Let $P$ be the orthogonal projection onto $W^\perp$, so that
$\operatorname{tr}P\geq d-1$.  Taking the trace after multiplication by
$P$ preserves the stated positive-semidefinite inequality: indeed,
$\operatorname{tr}(PQ)\geq0$ for positive-semidefinite $P,Q$.
Writing $E_A(\mathcal P)$ for the edges of $A$ crossing the partition,
we obtain
\begin{align*}
 c(d-1)
 &\leq c\operatorname{tr}P
 \leq\operatorname{tr}\left(P\sum_{e\in A}z_ez_e^{\mathsf T}\right)\\
 &=\sum_{e\in E_A(\mathcal P)}\|Pz_e\|^2
 \leq\frac{1+\rho}{q}|E_A(\mathcal P)|.
\end{align*}
The within-part terms vanish by the choice of $P$.  Every partition
therefore has at least $\lfloor cq/(1+\rho)\rfloor(d-1)$ crossing
edges.  Theorem~\ref{thm:packing} completes the proof; if the floor is
zero, the conclusion is immediate.
\end{proof}

\section{One sparsification step}
\label{sec:sparsification}

We first derive the exact consequence of the Marcus--Spielman--Srivastava
random-vector theorem~\cite[Theorem~1.4]{MSS}
needed here.  The use of deterministic padding permits singular
covariance matrices and keeps both halves under simultaneous control.

\begin{lemma}[Two-sided splitting]
\label{lemma:two-system-split}
Let $w_e\in\mathbb R^D$ be finitely many vectors, let
$\Sigma=\sum_e w_ew_e^{\mathsf T}\preceq2I_D$, and suppose that
$\|w_e\|^2\leq\epsilon$ for every $e$, where $\epsilon>0$.
There is a subset $A$ of the indices for which
\begin{equation}
 \frac12\Sigma-\tau I_D
 \preceq\sum_{e\in A}w_ew_e^{\mathsf T}
 \preceq\frac12\Sigma+\tau I_D,
 \qquad \tau=2\sqrt\epsilon+\epsilon.
 \label{eq:two-sided-split}
\end{equation}
\end{lemma}

\begin{proof}
Independently for each $e$, let $v_e\in\mathbb R^{2D}$ equal
$(w_e,0)$ or $(0,w_e)$, each with probability $1/2$.  Then
\[
 M:=\sum_e\mathbb E[v_ev_e^{\mathsf T}]
   =\operatorname{diag}(\Sigma/2,\Sigma/2)\preceq I_{2D},
 \qquad \mathbb E\|v_e\|^2=\|w_e\|^2\leq\epsilon.
\]
To make the covariance exactly isotropic, diagonalize the
positive-semidefinite matrix $I_{2D}-M$.  For each positive eigenvalue
$\lambda$ with unit eigenvector $f$, put
$m_\lambda=\lceil\lambda/\epsilon\rceil$ and introduce
$m_\lambda$ deterministic vectors equal to
$\sqrt{\lambda/m_\lambda}\,f$.
There are finitely many such vectors; each has squared norm at most
$\epsilon$, and their outer products sum to $I_{2D}-M$.
Deterministic vectors have singleton support and are independent of
all other vectors.  Consequently Theorem~\ref{thm:mss} applies to the
enlarged family.

In its successful outcome, let $A$ be the indices assigned to the first
block, and set $S_A=\sum_{e\in A}w_ew_e^{\mathsf T}$.
Subtracting the fixed padding matrix from the theorem's conclusion
gives
\[
 \operatorname{diag}(S_A,\Sigma-S_A)
 \preceq M+\bigl((1+\sqrt\epsilon)^2-1\bigr)I_{2D}
 =M+\tau I_{2D}.
\]
The first block gives the upper bound in
\eqref{eq:two-sided-split}; the second block gives its lower bound.
\end{proof}

\begin{proposition}[Packing-preserving sparsification]
\label{prop:step}
There is an absolute integer $r_{\mathrm{step}}$ such that the following
holds for every integer $r\geq r_{\mathrm{step}}$.  If $H$ has at
least two vertices and contains $r$ edge-disjoint spanning trees, set
\[
 u=r^{-1/16},\qquad q=\lfloor(1-u)r\rfloor,
 \qquad \epsilon=6r^{-1/2},\qquad
 \tau=2\sqrt\epsilon+\epsilon.
\]
Then $H$ has a spanning subgraph $H'$ containing at least
\begin{equation}
 r'=\left\lfloor\frac{(1-4u)r}{2}\right\rfloor
 \label{eq:step-packing-count}
\end{equation}
edge-disjoint spanning trees, such that every nonempty proper
$S\subset V(H)$ satisfies
\begin{equation}
 |\delta_{H'}(S)|\leq
 \left(\frac12+3\tau\right)|\delta_H(S)|.
 \label{eq:step-cut-bound}
\end{equation}
The cutoff may be chosen so that, in addition,
\begin{equation}
 \frac5{16}\leq\frac{r'}r\leq\frac12,
 \qquad
 \frac{\frac12+3\tau}{r'/r}\leq1+15u.
 \label{eq:step-ratio-bound}
\end{equation}
\end{proposition}

\begin{proof}
Write $n=|V(H)|$. Choose the cutoff large enough for
Theorem~\ref{thm:extraction}, and
also require $r\geq16^{16}$, so $u\leq1/16$.
The extraction theorem supplies a positive definite matrix $X$ with
\[
 X\succeq L_H,
 \qquad
 \mathbf1_S^{\mathsf T}X\mathbf1_S
 \leq3\mathbf1_S^{\mathsf T}L_H\mathbf1_S
 \quad(\varnothing\ne S\ne V),
\]
and $q$ edge-disjoint spanning trees all of whose edges satisfy
$b_e^{\mathsf T}X^{-1}b_e\leq4r^{-1/2}$.
Let $B$ be precisely the union of these trees.  Apply
Lemma~\ref{lem:auxiliary-representation} to $B$ with $\rho=u$, and put
\[
 a_e=X^{-1/2}b_e,\qquad w_e=(a_e,z_e),
 \qquad\Sigma=\sum_{e\in B}w_ew_e^{\mathsf T}.
\]
The covariance includes cross blocks:
\[
 \Sigma=\begin{pmatrix}
 X^{-1/2}L_BX^{-1/2}&\displaystyle\sum_{e\in B}a_ez_e^{\mathsf T}\\
 \displaystyle\sum_{e\in B}z_ea_e^{\mathsf T}&I_{n-1}
 \end{pmatrix}.
\]
They need not vanish.  Nevertheless, for any vectors $x,y$ in the two
coordinate spaces,
\begin{align*}
 (x,y)^{\mathsf T}\Sigma(x,y)
 &=\sum_{e\in B}\bigl(\langle x,a_e\rangle+
                          \langle y,z_e\rangle\bigr)^2\\
 &\leq2x^{\mathsf T}X^{-1/2}L_BX^{-1/2}x+2\|y\|^2\\
 &\leq2(\|x\|^2+\|y\|^2),
\end{align*}
because $L_B\preceq L_H\preceq X$.  Hence
$0\preceq\Sigma\preceq2I$ without any invertibility assumption on
$\Sigma$.

Since $q\geq(1-u)r-1\geq r/2$, and $r\geq4$,
\[
 \frac{1+u}{q}\leq\frac{17}{8r}\leq2r^{-1/2}.
\]
Thus every $w_e$ has squared norm at most
$4r^{-1/2}+(1+u)/q\leq\epsilon$.
Choose the subset $A\subseteq E(B)$ supplied by
Lemma~\ref{lemma:two-system-split}.  The two principal blocks of the
\emph{same} inequality \eqref{eq:two-sided-split} give
\begin{equation}
 \sum_{e\in A}z_ez_e^{\mathsf T}\succeq(1/2-\tau)I_{n-1},
 \qquad
 X^{-1/2}L_AX^{-1/2}
 \preceq\tfrac12X^{-1/2}L_BX^{-1/2}+\tau I_n.
 \label{eq:same-subset-two-constraints}
\end{equation}
Congruence by $X^{1/2}$ in the second inequality yields
\[
 L_A\preceq\tfrac12L_B+\tau X
       \preceq\tfrac12L_H+\tau X.
\]
Testing on $\mathbf1_S$ proves \eqref{eq:step-cut-bound} for
$H'=(V,A)$, simultaneously for all cuts.

For the packing assertion, write the errors in terms of $u$:
\[
 \epsilon=6u^8,\qquad\tau=2\sqrt6\,u^4+6u^8\leq u/2
 \quad(0<u\leq1/16).
\]
For example, dividing by $u$ and using $\sqrt6<3$ bounds the left
ratio by $6(16^{-3}+16^{-7})<1/2$.
In particular $c=1/2-\tau>0$, so the first inequality in
\eqref{eq:same-subset-two-constraints} and
Lemma~\ref{lem:projected-trace-packing} give at least
$\lfloor(1/2-\tau)q/(1+u)\rfloor$ trees.
The rounding margin is explicit:
\begin{align*}
 &(1-2\tau)\bigl((1-u)r-1\bigr)
            -(1+u)(1-4u)r\\
 &\quad\geq(1-u)\bigl((1-u)r-1\bigr)
            -(1+u)(1-4u)r\\
 &\quad=r(u+5u^2)-(1-u)\geq0,
\end{align*}
where $ru=r^{15/16}\geq1$.  Consequently
\[
 \frac{(1/2-\tau)q}{1+u}\geq\frac{(1-4u)r}{2}.
\]
Taking floors proves \eqref{eq:step-packing-count}.

Finally, the floor in that formula and $1/r\leq u$ give
\[
 \frac{r'}r\geq\frac12-2u-\frac1r
              \geq\frac12-3u\geq\frac5{16},
 \qquad\frac{r'}r\leq\frac12.
\]
Using $\tau\leq u/2$ once more, we obtain
\[
 \frac{\frac12+3\tau}{r'/r}
 \leq\frac{1+3u}{1-6u}
 =1+\frac{9u}{1-6u}\leq1+15u,
\]
which proves \eqref{eq:step-ratio-bound}.
\end{proof}
